% New slide.

\begin{frame}[t, fragile]{Haskell functions that print to IO : putStrLn}

Haskell has a number of functions that can print text to IO. Some of these include;

\begin{minted}[linenos=false]{haskell}
    putStrLn
    print
    printf
\end{minted}

But apart from printing text to IO, what values do these functions return?

putStrLn has the following function type signature;

\begin{minted}[linenos=false]{haskell}
putStrLn :: String -> IO ()
\end{minted}

This means that this function takes a single argument of type String, performs an IO action using it -- i.e. prints the
value of the argument to IO, and then returns a value of type (). Recall from earlier, that () denotes the unit type.
According to ChatGPT, () carries no information. It is used in this case to indicate that the function has completed
its action successfully.

\end{frame}


% New slide.

\begin{frame}[t, fragile]{Haskell functions that print to IO : print}

What about the function print? Its function type signature is;

\begin{minted}[linenos=false]{haskell}
print :: Show a => a -> IO ()
\end{minted}

But what does Show denote in this function's type signature? In Haskell, Show is what's known as a type class. That is,
a class (or collection) of types -- which in this case, can be converted into a human-readable string and ``shown'' on
IO.

Recall that we saw this type of syntax earlier when we declared/defined our parametrically polymorphic version of the
function addTwoNumbers; where instead of class type Show, it utilised class type Num.

\begin{minted}[linenos=false]{haskell}
addTwoNumbers :: Num a => a -> a -> a
\end{minted}

Some examples of types which belong to the Sum type class are;

\begin{minted}[linenos=false]{haskell}
Char
String
Int
Double
Float
\end{minted}

\end{frame}


% New slide.

\begin{frame}[t, fragile]{Haskell functions that print to IO : print}

That is, all of the native Haskell types which were just listed, i.e. Char, String, Int, Double, and Float, are all
instances of Show.

Haskell class types are a bit like C++ abstract base classes (ABCs). Consider the declaration of a very basic C++ ABC
which is shown below.

\begin{minted}[linenos=true]{haskell}
class
Showable
{
public :
    virtual
    std::string
    show()
    const = 0;
};
\end{minted}

Since this class is declared as being ``pure virtual'' -- as denoted by the presence of the ``virtual'' and
``const = 0'' tags, it means that it cannot be instantiated. This means that we cannot create objects of it. Haskell
type classes are exactly the same; thus we cannot instantiate type class Sum or Num in Haskell.

\end{frame}


\begin{frame}[t, fragile]{Haskell functions that print to IO : printf}

Last of all, let us consider the function ``printf''. It has the following function type signature;

\begin{minted}[linenos=false]{haskell}
printf :: PrintfType r => String -> r
\end{minted}

printf is an interesting function, as what it returns can vary, depending on what is expected of it. The following
invocation of the function will return a value of IO ();

\begin{minted}[linenos=false]{haskell}
printf "The answer is %d\n" 42
\end{minted}

whilst the following invocation of the function will return a value of String, as that is what the surrounding code
is expecting;

\begin{minted}[linenos=false]{haskell}
s :: String
s = printf "Pi = %.3f" pi
\end{minted}

Using printf in this way in Haskell appears to yield similar behaviour to that of stringstreams in C++. Consider the
following snippet of C++ code. To help try and convey this idea, consider the following C++ code snippet;

\begin{minted}[linenos=false]{C++}
#include <sstream>

std::stringstream ss;

ss << std::fixed << std::setprecision(3) << "Pi = " << std::numbers::pi;
\end{minted}

\end{frame}


\begin{frame}[t, fragile]{Haskell functions that print to IO : printf}

In the Haskell example, the string is constructed by the printf function and then bound to the value s.

In the C++ example, the string is constructed by the stream insertion operators (<<) and then the result is stored into
stringstream variable ss.

\end{frame}


\begin{frame}[t, fragile]{Chaining Haskell expressions together and the do keyword.}

Remember earlier, that we stated; \\

\begin{quote}
Haskell can be thought of as being an expression-oriented language. Every expression has a value, and a function body is
an expression.
\end{quote}

Recall that our function ``addTwoNumbers\_printAndReturnResult'' from earlier, made use of the ``do'' keyword. But what
would this function look like if we didn't use the ``do'' keyword? The answer is shown in the following code snippet;

\begin{minted}[linenos=true]{haskell}
addTwoNumbers_printAndReturnResult :: Int -> Int -> IO Int
addTwoNumbers_printAndReturnResult arg1 arg2 =

    let

        sum = arg1 + arg2

    in

        printf "Sum = %d" sum >>= \_ ->
        return sum
\end{minted}

\end{frame}


\begin{frame}[t, fragile]{Chaining Haskell expressions together and the do keyword.}

The Use \mintinline{cpp}|>>=| operator in the previous code snippet says to take the result of the monadic
\mintinline{cpp}|printf| function and pass it on to the rest of the block of expressions that form the body of the
function.

The \mintinline{cpp}|>>=| operator is called the bind operator. According to ChatGPT, it is one of the core operations
of the Monad type class. Furthermore, its type signature is as follows;

\begin{minted}[linenos=false]{haskell}
(>>=) :: Monad m => m a -> (a -> m b) -> m b
\end{minted}

\end{frame}
